
\documentclass{article}
\usepackage{colortbl}
\usepackage{makecell}
\usepackage{multirow}
\usepackage{supertabular}

\begin{document}

\newcounter{utterance}

\centering \large Interaction Transcript for game `ifeval', experiment `ifeval', episode 59 with Qwen3.5{-}9B.
\vspace{24pt}

{ \footnotesize  \setcounter{utterance}{1}
\setlength{\tabcolsep}{0pt}
\begin{supertabular}{c@{$\;$}|p{.15\linewidth}@{}p{.15\linewidth}p{.15\linewidth}p{.15\linewidth}p{.15\linewidth}p{.15\linewidth}}
    \# & \multicolumn{2}{c}{Player} && \multicolumn{2}{c}{Game Master} \\
    \hline

    \theutterance \stepcounter{utterance}  
    & & & \multicolumn{4}{p{0.6\linewidth}}{
        \cellcolor[rgb]{0.9,0.9,0.9}{
            \makecell[{{p{\linewidth}}}]{
                \texttt{\tiny{[P1$\langle$GM]}}
                \texttt{Write a short proposal for a new research project that investigates how language evolves over time. I want to make it challenging, so:} \\
\texttt{1. Do not include any commas in your response.} \\
\texttt{2. Do not include the letter "c" anywhere in your response.} \\
\texttt{3. Your response should contain at least 250 words.} \\
            }
        }
    }
     \\ \\

    \theutterance \stepcounter{utterance}  
    & \multicolumn{4}{p{0.6\linewidth}}{
        \cellcolor[rgb]{0.561,0.737,0.561}{
            \makecell[{{p{\linewidth}}}]{
                \texttt{\tiny{[P1$\rangle$GM]}}
                \texttt{This ambitious study aims to map the shifting patterns of human speech across many decades without relying on traditional written records. We will analyze vast digital archives containing millions of audio clips from diverse regions around the globe. Our team will use advanced artificial intelligence tools to detect subtle changes in pronunciation and grammar structures over time. The goal is to understand how new words enter our daily vocabulary while old terms slowly fade away from common use. We will focus on how social media platforms accelerate these linguistic transformations compared to slower historical trends. Researchers will examine how slang spreads rapidly through online networks before becoming standard usage in formal settings. This project requires collecting data from hundreds of different communities to ensure broad representation across all cultures. We must build custom software capable of processing huge amounts of unstructured sound files efficiently. The team will develop new algorithms that can identify emerging dialects before they become widely recognized by scholars. Our findings could reveal how technology reshapes the way people communicate with one another every single day. We plan to publish results showing clear evidence of rapid change in modern language systems. This work will challenge existing theories about how slowly languages typically evolve over centuries. By studying current trends we might predict future developments in global communication styles. The project demands close collaboration between linguists computer scientists and sociologists working together daily. We will hold weekly meetings to review progress and adjust our methods as needed for better accuracy. Funding will support travel for fieldwork in remote areas where digital data is scarce or missing entirely. Our ultimate aim is to create a comprehensive database tracking every major shift in spoken English since the year two thousand. This resource will help educators understand how students learn new words in today's fast paced world. We believe this research will open new doors for understanding human connection through shared verbal expressions. The study will also explore how migration influences local dialects when people move between countries frequently. Our approach avoids using static snapshots in favor of continuous monitoring of living language systems. This dynamic perspective offers unique insights into the fluid nature of human expression today. We expect to uncover surprising patterns that contradict previous assumptions about language stability. The final report will serve as a foundational text for future generations studying this fascinating subject matter. Our team remains committed to delivering high quality findings that advance the field significantly. This endeavor represents a critical step forward in our collective understanding of human communication evolution. We look forward to sharing these discoveries with the wider academic community soon. The project timeline spans five years allowing ample time for thorough analysis and verification of all data points collected. Every member of the group brings unique skills essential for completing this complex and rewarding task successfully.} \\
            }
        }
    }
    & & \\ \\

    \theutterance \stepcounter{utterance}  
    & & & \multicolumn{2}{p{0.3\linewidth}}{
        \cellcolor[rgb]{0.9,0.9,0.9}{
            \makecell[{{p{\linewidth}}}]{
                \texttt{\tiny{[GM$|$GM]}}
                {'punctuation:no_comma': {}, 'keywords:letter_frequency': {'let_relation': 'less than', 'letter': 'c', 'let_frequency': 1}, 'length_constraints:number_words': {'relation': 'at least', 'num_words': 250}}
            }
        }
    }
    & & \\ \\

    \theutterance \stepcounter{utterance}  
    & & & \multicolumn{2}{p{0.3\linewidth}}{
        \cellcolor[rgb]{0.9,0.9,0.9}{
            \makecell[{{p{\linewidth}}}]{
                \texttt{\tiny{[GM$|$GM]}}
                \texttt{game\_result = LOSE} \\
            }
        }
    }
    & & \\ \\

\end{supertabular}
}

\end{document}
